\subsection{Smaller witnesses and complete-call measurements}

The final measurements run serially, without tests, audits, builds or commits
overlapping them. Each case has one warm-up round and five measured rounds.
Each round randomly orders old, old A/A, new and new A/A calls. The tables
give each primary arm's median elapsed time and the median of paired ratios;
the latter need not equal a ratio of the two medians. All 200 measured and
40 warm-up calls complete. Baseline and new certificate hashes are stable
within their own implementations; their coordinate hashes agree across
implementations. Different source-proof schemas intentionally produce
different complete certificate hashes.

The first experiment times source-bound independent verification of a
precomputed positive certificate, together with source parsing, certificate
hashing and serialized-byte counting. Producing the certificates is outside
the timed sample. It includes all geometry and cohomology checks needed by
each respective verifier, with no discovery cache passed into replay.
The mode's total elapsed time, including preparation, is 11.136 seconds.

\begin{center}
\small\begin{tabular}{lrrrrr}
\toprule
Input & Old (ms) & New (ms) & Old/new & Old A/A & New A/A \\
\midrule
circle-8 & 55.735 & 21.468 & 2.591 & 0.993 & 1.006 \\
circle-16 & 121.078 & 43.413 & 2.786 & 0.990 & 0.998 \\
circle-32 & 273.556 & 92.205 & 2.982 & 0.973 & 0.992 \\
Optimized positive & 30.206 & 15.744 & 1.937 & 1.001 & 1.000 \\
\bottomrule
\end{tabular}

\end{center}

Complete replay improves $1.94$--$2.98\times$ on these four cases. The
canonical compact JSON witnesses for the 8-, 16- and 32-crossing circles
shrink from 112,472, 237,105 and 491,293 bytes to 30,077, 60,649 and
121,929 bytes. The optimized positive shrinks from 51,195 to 21,671
bytes, including its height data and dual matching. These sizes include
the common source triangulation and PD; they are not sizes of an isolated
connectivity subcertificate.

The second experiment starts with a source PD and runs the complete
recognizer, including source construction, cohomology discovery, an optional
span search, the stage's internal checks, the exact fallback when required,
and an additional independent replay of any returned native proof. It also
hashes and counts the serialized certificate. Earlier shortcuts are
deliberately disabled so that every case reaches the native stage. Both
arms have no local seed work cap and retain a twelve-second global deadline.
All calls complete before that deadline; the mode takes 53.089 seconds.

\begin{center}
\small\begin{tabular}{lrrrrr}
\toprule
Input & Old (ms) & New (ms) & Old/new & Old A/A & New A/A \\
\midrule
circle-8 & 260.410 & 64.166 & 4.004 & 1.001 & 0.921 \\
circle-16 & 744.828 & 131.303 & 5.756 & 1.002 & 0.998 \\
circle-32 & 2345.215 & 269.486 & 8.702 & 1.005 & 1.002 \\
Optimized positive & 233.967 & 114.397 & 2.050 & 0.999 & 0.999 \\
Genus-one miss & 72.198 & 49.426 & 1.467 & 1.006 & 0.992 \\
Trefoil & 71.814 & 49.088 & 1.459 & 1.001 & 1.008 \\
\bottomrule
\end{tabular}

\end{center}

The four native positives improve $2.05$--$8.70\times$. The 32-crossing
circle's median falls from 2.345 seconds to 0.269 seconds. The genus-one
unknot miss and trefoil still use the exact Khovanov fallback, returning
\code{UNKNOT} and \code{KNOTTED} respectively; their complete calls
improve about $1.46\times$ by making the preceding unsuccessful stage
cheaper. The new-arm A/A ratio on the smallest circle is 0.921, indicating
more variation there than on the other cases; its paired old/new gain is
still about four. These selected-workload gains establish neither ordinary
default-portfolio acceleration nor a general complexity theorem. The
default-budget source audit above separately establishes the newly
completed 32-crossing positive.
